\begin{table}[!htbp]
\centering
\caption{Efeito do filtro nas 19 STNs suplementares agregadas.}
\label{tab:res-stn-representacoes}
\begin{singlespace}
\normalsize
\setlength{\tabcolsep}{5pt}
\renewcommand{\arraystretch}{1.15}
\begin{tabular}{lrrrr}
\hline\hline
\textbf{Representação} & \textbf{Registros} & \textbf{Nós} & \textbf{Arestas} & \textbf{\shortstack{Transições\\ponderadas}} \\
\hline
Bruta & 36.593 & 2.682 & 3.874 & 34.310 \\
Filtrada & 14.942 & 744 & 830 & 4.391 \\
\hline\hline
\end{tabular}
\par\smallskip
{\small Nota: somas das 19 redes, com nós distintos dentro de cada instância. A rede bruta preserva as transições observadas; a filtrada pode conectar estados não adjacentes na sequência original.\par}
{\fonte{elaboração própria a partir dos dados experimentais.}}
\end{singlespace}
\end{table}
